% Revisione: stati di animazione della wall slide e layer d'animazione disaccoppiati, 25 ottobre 2026.
% Compilare con: pdflatex wall-anim-2026-10-25.tex (due volte, per l'indice)
\documentclass[11pt,a4paper]{article}
\usepackage[utf8]{inputenc}
\usepackage[T1]{fontenc}
\usepackage[italian]{babel}
\usepackage{lmodern}
\usepackage[margin=2.3cm]{geometry}
\usepackage{listings}
\usepackage{xcolor}
\usepackage{booktabs}
\usepackage{hyperref}
\usepackage{enumitem}

\lstset{basicstyle=\ttfamily\small,breaklines=true,frame=single,columns=fullflexible,
  keywordstyle=\bfseries,commentstyle=\itshape\color{gray},
  literate={à}{{\`a}}1 {è}{{\`e}}1 {é}{{\'e}}1 {ì}{{\`i}}1 {ò}{{\`o}}1 {ù}{{\`u}}1}
\newcommand{\code}[1]{\texttt{\detokenize{#1}}}

\title{Animare la wall slide\\\large Stati di animazione e layer disaccoppiati dall'ABP}
\author{Ramo \code{claude/project-thread-2vuw3f} (da \code{claude/project-thread-v1plhd})}
\date{25 ottobre 2026}

\begin{document}
\sloppy
\maketitle
\tableofcontents

\section{Cosa c'è in questo documento}
La wall slide v2 (\code{WallSlidingV2}, \code{WallClimbingV2}) ora espone all'animazione tutto ciò che serve per più stati:
entrata, uscita, scivolata da fermo, verso sinistra, verso destra e skid sopra una certa velocità. In più c'è il primo pezzo
dello schema generale che hai chiesto: le funzionalità portano le loro animazioni in \emph{layer} separati, scelti da dati,
e l'ABP principale non le nomina mai. Le varianti per skin (con hitbox identiche) sono una proposta, in attesa delle tue risposte
(sezione~\ref{sec:domande}).

Il codice è scritto nel cloud e \textbf{non è ancora stato compilato}: va compilato sul fisso (sezione~\ref{sec:controlli}).
Le modifiche all'ABP non sono fatte: l'ABP è un Blueprint condiviso e l'editor del fisso è in uso dall'altro thread.

\section{Stati della wall slide}
\subsection{Ragionamento}
L'animazione non fa parte della simulazione: si ricava a ogni frame, sul game thread, dallo stato di presentazione di Mover
(lo stato sincronizzato già smussato sui proxy). Per questo può tenere i suoi timer (durata dell'entrata e dell'uscita) senza
toccare predizione e rollback, e le modalità di movimento non sanno nulla dell'animazione.

``Sul muro'' significa: la modalità attiva è della classe \code{UGSTN_WallTraverseModeBase} (slide e climb v2). Così il passaggio
slide $\leftrightarrow$ climb sul posto non è un'uscita seguita da un'entrata. Le modalità legacy sono ignorate, come hai chiesto,
ma basta cambiare la classe nelle impostazioni per includerle.

La velocità viene scomposta sul piano del muro: \emph{laterale} (positiva verso la destra del personaggio, che guarda il muro),
\emph{verticale} (lungo la gravità proiettata sul muro, positiva in su; funziona con qualsiasi direzione della gravità) e
\emph{velocità sul piano}. Lo skid è una regola di velocità perché il freno non è nello stato sincronizzato: un'entrata veloce
frena e passa attraverso lo skid in modo naturale. Ogni soglia ha un'isteresi, così lo stato non sfarfalla al confine.

\subsection{Pseudocodice}
\begin{lstlisting}
ogni frame (game thread, prima del grafo):
  ingressi = modalità attiva, velocità, normale filtrata del muro, gravità
  fase:
    None/Exit e sul muro       -> Enter (o Active se EnterDuration = 0)
    Enter da EnterDuration     -> Active
    Enter/Active e non sul muro -> Exit (salva la modalità che subentra)
    Exit da ExitDuration       -> None
  se sul muro:
    (laterale, verticale, velocità) = scomponi(velocità, normale, gravità)
    se non climb e velocità >= (skid prima ? 650 : 750)      -> Skid
    altrimenti se |laterale| >= (stesso lato prima ? 80 : 120) -> SlideLeft / SlideRight
    altrimenti                                               -> Idle
  fuori dal muro: stato e velocità restano gli ultimi (l'uscita può sceglierne la variante)
\end{lstlisting}

\subsection{Classi e strumenti Unreal}
\begin{itemize}[leftmargin=*]
\item \code{FGSTN_WallAnimSettings}, \code{FGSTN_WallAnimState}, \code{FGSTN_WallAnimInputs}, enum \code{EGSTN_WallAnimPhase}
  e \code{EGSTN_WallSlideAnimState} (\code{WallMovement/Animation/GSTN_WallAnimTypes.h}).
\item \code{UGSTN_WallAnimUtils} (\code{UBlueprintFunctionLibrary}, senza stato): \code{SplitWallVelocity}, \code{ClassifySlide},
  \code{StepWallAnimState}, \code{GatherInputs} (legge \code{UMoverComponent::GetMovementMode}, \code{GetVelocity},
  \code{GetUpDirection} e \code{FGSTN_WallMovementSyncState::WallNormal}).
\item \code{UGSTN_WallAnimInstance} (\code{UAnimInstance}): genitore C++ del layer della wall slide. In
  \code{NativeUpdateAnimation} aggiorna \code{WallAnim}; il grafo lo legge con \emph{Property Access} (thread safe).
\item Test \code{GSTN.Wall.Anim.*}: scomposizione con gravità diverse, classificazione con isteresi, passi di fase
  (entrata, slide$\to$climb senza uscita, ripresa durante l'uscita, durate a zero).
\end{itemize}

\subsection{Il grafo del layer (con animazioni segnaposto)}
\begin{lstlisting}
ABP_GSTN_WallLayer (genitore GSTN_WallAnimInstance, implementa ALI_GSTN_Wall), layer WallLayer:
  Enter --(Phase == Active)--> Loop
  Loop = macchina interna su SlideState: Idle / SlideLeft / SlideRight / Skid
         (sequenze o blend space su LateralAlpha, VerticalAlpha)
  Loop o Enter --(Phase == Exit)--> Exit
  Exit --(Phase == Enter)--> Enter   (ripresa del muro durante l'uscita)
\end{lstlisting}

\section{Animazioni per funzionalità, disaccoppiate dall'ABP}\label{sec:layer}
\subsection{Ragionamento}
Hai chiesto che wall slide, spin, trapano e simili portino le proprie animazioni senza legarle all'ABP, e che una skin (o
qualche sua caratteristica) possa cambiarle lasciando identica la hitbox. È lo schema di Lyra: ogni funzionalità ha un
\emph{linked anim layer} (un anim Blueprint a parte, con la sua interfaccia), collegato con \code{LinkAnimClassLayers};
una variante è un anim Blueprint figlio. In più la scelta avviene per tag, per priorità tra le sorgenti.

\subsection{Parte A (fatta)}
\begin{itemize}[leftmargin=*]
\item \code{IGSTN_OverrideLayerSource}: un layer dice se vuole prendere la posa (\code{IsOverrideLayerActive}), con quale
  tempo di blend e con che priorità.
\item \code{UGSTN_CharacterAnimInstance}: genitore C++ dell'ABP principale. Non conosce nessuna funzionalità: a ogni frame chiede
  ai layer collegati (\code{GetLinkedAnimInstances}) chi è attivo ed espone \code{bOverrideLayerActive} e
  \code{OverrideBlendTime}. I layer collegati si aggiornano prima dell'istanza principale, quindi la risposta è del frame corrente.
\item Nell'ABP principale servono solo: il reparent, un \emph{Blend Poses by bool} e un nodo \emph{Linked Anim Layer} per ogni
  interfaccia di layer. Nessuna modifica per ogni nuova variante.
\end{itemize}

\subsection{Parte B (proposta)}
\begin{lstlisting}
plugin GSTNAnimation, modulo GSTNAnimationCore (puro: Engine, GameplayTags)
UGSTN_AnimSetDefinition (data asset): SetTag (Anim.Set.Wall/Spin/Drill), LayerClass,
                                      MontageOverrides { Reference, Visual }
FGSTN_AnimSetRule: Set, When (query sui tag cosmetici: SkinTags + Anim.Trait.*), Priority
UGSTN_AnimSetCollection: Rules           (su personaggio, kit, arma, skin)
UGSTN_AnimSetComponent (solo presentazione):
  a ogni cambio di sorgente o di tag (mai a ogni frame):
    per ogni SetTag: vincitore = max(priorità sorgente, priorità regola) tra le regole che combaciano
    carica in asincrono, LinkAnimClassLayers(LayerClass), mappa Reference -> Visual
  ResolveVisualMontage(Reference)
\end{lstlisting}

\subsection{La hitbox non cambia: le regole}
\begin{enumerate}[leftmargin=*]
\item La simulazione (timeline, finestre di hit e hurtbox, \code{BakedSocket}, mosse in coda di Mover) legge solo il montage
  \emph{di riferimento} e non dipende dal componente degli anim set.
\item I socket si fanno bake solo dal montage di riferimento; una variante visiva non si fa mai bake.
\item Un validatore dell'editor richiede che ogni variante abbia le stesse sezioni, le stesse durate (entro un frame), lo stesso
  play rate e nessun root motion diverso.
\item Le hurtbox seguono l'\code{HurtboxSet} del kit, non la skin.
\end{enumerate}

\section{Domande per te}\label{sec:domande}
\begin{enumerate}[leftmargin=*]
\item Nuovo plugin \code{GSTNAnimation} (consigliato) o tipi nel modulo del progetto?
\item Varianti per skin con regole a tag nelle collezioni (consigliato: le skin restano invariate, le caratteristiche sono tag)
  o una lista diretta di anim set su ogni skin?
\item Le varianti devono avere le stesse sezioni e durate del riferimento (consigliato) o si permettono durate diverse con
  riallineamento ai frame del riferimento?
\end{enumerate}

\section{Cosa controllare a mano}\label{sec:controlli}
\begin{enumerate}[leftmargin=*]
\item \textbf{Compilazione}: build completa del modulo (file nuovi e nuovi \code{UCLASS}/\code{USTRUCT}), poi Session Frontend
  \code{GSTN.Wall.Anim} (3 test).
\item \textbf{Editor}, quando è libero (Blueprint condiviso, faccio io o tu):
  \begin{enumerate}
  \item creare \code{ALI_GSTN_Wall} con il layer \code{WallLayer};
  \item creare \code{ABP_GSTN_WallLayer} (genitore \code{GSTN_WallAnimInstance}) con il grafo della sezione 2.4 e clip GASP
    segnaposto;
  \item fare il reparent di \code{SandboxCharacter_Mover_ABP} a \code{GSTN_CharacterAnimInstance}, aggiungere il blend e il
    nodo \code{WallLayer} con istanza \code{ABP_GSTN_WallLayer}.
  \end{enumerate}
\item \textbf{PIE}, mappa \code{L_WallSlideAB}, con \code{showdebug animation} o il debugger dell'ABP sul layer:
  \begin{enumerate}
  \item salto contro il muro da fermo: Enter, poi Active/Idle;
  \item stick a destra e a sinistra: SlideRight / SlideLeft, nessuno sfarfallio vicino a 80--120\,cm/s;
  \item entrata in corsa veloce: Skid, che finisce sotto 650\,cm/s;
  \item passaggio su una superficie da climb: resta Active (nessuna uscita), \code{bIsClimbing} vero, niente skid;
  \item salto dal muro o atterraggio: Exit con \code{ExitModeName} giusto, poi None;
  \item in listen server, un client che guarda un altro giocatore: stessi stati (lo stato viene dalla presentazione smussata).
  \end{enumerate}
\end{enumerate}

\end{document}
